\section{Introduction}
%Content moderation has become a cornerstone of trust and safety on large social platforms.

% maybe start with human-ai worksflows and later give content moderation as an exmaple.
%Uncertainty estimation has become central to deploying Large Language Models (LLMs) in high-stakes human-AI workflows, such as content moderation.  Moderation systems aim to achieve a balance between the scale and cost-efficiency of AI, and the superior accuracy of human experts, contributing to safer online environments. The rapid growth of user-generated content makes manual review infeasible, thereby motivating the growing deployment of Large Language Models (LLMs), for scalable moderation. When LLMs are used in human-in-the-loop  workflows, their utility depends not only on their accuracy, but also on their ability to signal when they are likely to be wrong.  The effectiveness of such pipelines hinges on a single decision that repeats millions of times per day: should the system trust the LLM judgment or escalate to a human reviewer?
Human–AI collaboration is at the heart of modern content moderation, which serves as a cornerstone of online trust and safety. Effective moderation is essential for protecting individuals from harmful or misleading content, and the pursuit of scalable and responsible approaches contributes to a safer online environment. To meet this need, moderation systems must balance the nuanced accuracy of human experts with the scalability and cost-efficiency of AI. The rapid growth of user-generated content has made reliance solely on manual review infeasible, motivating the growing deployment of Large Language Models (LLMs) into moderation workflows. When LLMs are used in human-in-the-loop workflows, their utility depends not only on their accuracy, but also on their ability to signal when they are likely to make an incorrect decision.  The effectiveness of such pipelines depends on a single decision that repeats billions of times per day: should the system trust the LLM's judgment or escalate to a human reviewer?

Accordingly, we focus on two research questions:\\
%This question drives our work and leads to two guiding research questions:\\
\noindent \firstmention{RQ1:} Can we accurately predict when LLMs are likely to be incorrect in moderation tasks? \\
\noindent \firstmention{RQ2:} Can selective escalation based on these predictions reduce the overall cost in human-AI moderation workflows?

This trust-or-escalate decision is, at its core, an uncertainty estimation problem. 
Recent work~\cite{shorinwa2025surveyuncertaintyquantificationlarge,  farr+al:25, kapoor2024large, cohen2024i, efficient-and-effective, liu+al:24} has explored a variety of signals for estimating LLM uncertainty. Approaches range from token-level entropy and log-probability thresholds~\cite{jiang+al:21, huly+al:25} to ensemble-based sampling~\cite{farr+al:25} and explicit self-expressed confidence~\cite{tian+al:23, kadavath+al:22, lin+al:22, xiong+al:23}.

In the realm of content moderation, current LLM uncertainty metrics fail to capture the economic value of moderation failure. To address this, we formulate the decision as a cost minimization problem, balancing between the cost of misclassification and the cost of human review. Moreover, previous work~\cite{shorinwa2025surveyuncertaintyquantificationlarge,jiang+al:21} showed that no single universal uncertainty metric has been proven to be reliable across tasks and models. These findings motivate the pursuit of a robust approach that can perform reliably at scale in high-stakes domains such as content moderation. 

To this end, we introduce a novel framework that builds on this foundation, enabling a unified approach to selective escalation. Unlike prior supervised uncertainty estimation approaches that typically rely on single or homogeneous signals~\cite{liu+al:24, jiang+al:21, hendrycks2017a}, or learning-to-defer methods that optimize deferral without uncertainty attribution~\cite{geifman2017selectiveclassificationdeepneural, mozannar2021consistentestimatorslearningdefer}, we fuse diverse gray- and black-box output signals and incorporate moderation-specific abstention with attribution, producing more comprehensive and explainable LLM Performance Predictors (LPPs). The approach is inspired by Query Performance Prediction (QPP)~\cite{carmel+Elad:10,hauff+al:08,shtok+al:16,shtok+al:12,tao+al:14,zhou+croft:07} in Information Retrieval, which studies how systems can learn to anticipate the success of individual queries. QPPs are commonly divided into pre-retrieval predictors and post-retrieval predictors. In this work, we focus on post-generation predictors, which are more scarce in the literature~\cite{huly+al:25}, due to strong empirical predictive power, as they leverage richer evidence from the actual model outputs and broad applicability for both open and closed models.

Uncertainty estimators for LLMs are often categorized by the level of access to model internals: white-box approaches leverage hidden states and gradients, gray-box approaches use output-side signals such as log-probabilities, and black-box approaches rely only on final predictions or agreement across multiple generations. In this work, we focus on gray- and black-box features, as they strike a practical balance between richness of signal, scale, and applicability across both off-the-shelf and open-source LLMs.

Beyond estimating the level of uncertainty, moderation teams also need to know the reason for the uncertainty. We therefore introduce novel moderation-oriented uncertainty attribution indicators that distinguish evidence deficits (e.g. missing transcript context, visually ambiguous frames, cross-lingual inconsistencies) from policy gaps (e.g., edge cases not covered by moderation guidelines, conflicting definitions across regions).
This enables practical routing: “Tough Calls” (aleatoric) are routed to senior reviewers, while “Policy Gaps” (epistemic) trigger policy updates or model retraining with active learning. In doing so, LPPs do not merely indicate uncertainty; they also shorten the policy-improvement loop and focus scarce human effort, thereby minimizing the cost of the entire moderation system.

We evaluate the LPPs across multiple families of LLMs, off-the-shelf (Gemini, GPT) and open-source (Llama, Qwen), on two moderation datasets, featuring multimodal and multilingual content across multiple risk categories: hate, violence, and more. Across settings, LPPs consistently improve the accuracy and cost frontiers relative to each existing uncertainty estimator. We release reproducibility code alongside the paper to enable practitioners to build upon this framework and facilitate future research on LLM Performance Prediction.\footnote{\url{https://github.com/ZEFR-INC/lpp-research}}

\myparagraph{Contributions} 
(1) We propose a novel framework for estimating LLM uncertainty, learning a meta-model based on LLM Performance Predictors (LPPs) derived from the LLM's output; (2) we demonstrate the merits of our method for selective escalation in a real-world hybrid human-AI workflow for content moderation compared to existing uncertainty estimators; and (3) we introduce novel moderation-oriented uncertainty attribution indicators and show their contribution to explainability, offering insights into when and why LLMs fail, thereby improving learning efficiency in human-AI workflows.
% The key contributions of this work can be summarized as follows:
% \begin{enumerate}
%     \item We propose a novel framework for estimating LLM uncertainty, learning a meta-model based on LLM Performance Predictors (LPPs) derived from the LLM's output.
%     \item  We demonstrate the merits of our method for selective escalation in a real-world hybrid human-AI workflow for content moderation - significantly improving accuracy and cost across state-of-the-art LLMs, compared to existing uncertainty estimators.
%     \item  We introduce novel moderation-oriented uncertainty attribution indicators and show their contribution to explainability, offering insights into when and why LLMs fail, thereby improving learning efficiency in human-AI workflows.
% \end{enumerate}




